%! Piecewise-Defined Functions — Unit 2, Lesson 2.6 — Guided Notes (Answer Key)
\documentclass[10pt]{article}
\usepackage{atda-article}
\usepackage{atda-key}   % pulls in -boxes; do NOT also load -boxes
\newcommand{\TallMath}[1]{$\displaystyle #1\rule[-1.4em]{0pt}{3.2em}$}
\setlength{\workrowsep}{6pt}

\begin{document}

\pageheader{Unit 2, Lesson 2.6}{Guided Notes --- Answer Key}
% NAMESTRIP: no \namedateperiod here — mirrors the blank. Teacher-only prose goes
% in the lesson plan as \begin{teachernote}[Guided Notes], never in a key.

\vspace{0.15in}

\begin{objectivebox}
\small
By the end of this lesson, I will be able to\dots
\begin{enumerate}[leftmargin=1.5em, itemsep=2pt, topsep=3pt]
  \item read a graph that takes \textbf{more than one rule}, and evaluate it by first asking
        \emph{which piece owns this input} and then using that piece's rule and only that one;
  \item explain why a piecewise-defined function is still \textbf{one} function --- every input in
        the domain belongs to exactly one piece, so every input still has exactly one output;
  \item use the \textbf{open and closed circles} at a boundary to say what the function's value is
        there, and to write a range with a bracket or a parenthesis;
  \item decide whether a graph is \textbf{continuous} at a boundary, and tell a \emph{corner} from a
        \emph{break};
  \item find all of Unit 2's characteristics on a piecewise graph --- domain and range, zeros and
        intercepts, increasing/decreasing/constant intervals, absolute maximum and minimum, end
        behavior, and asymptotes.
\end{enumerate}
\end{objectivebox}

\vspace{0.10in}

\boxguard
\begin{vocabbox}
\small
Fill in each term as we build it together. Every one of them is read straight off a graph.
\vocabans{Piecewise-defined function}{one function described by more than one rule, where each rule
    applies only on its own stretch of the domain. The stretches do not overlap, so every input still
    has exactly one output.}
\vocabans{Piece}{one of those rules together with the stretch of inputs it owns. ``Which piece owns
    this input?'' is the first question you ask before evaluating.}
\vocabans{Boundary point}{an input where one piece hands off to the next --- $m=3$ on the delivery
    fee, $h=40$ on the paycheck. It is the input to check whenever anything looks strange.}
\vocabans{Open circle / closed circle}{the marks at the end of a piece. A \emph{closed} (filled)
    circle means that endpoint is part of the graph; an \emph{open} circle means it is not. Exactly
    one filled circle sits above each input of the domain.}
\vocabans{Continuous / discontinuous}{a graph is \emph{continuous} on an interval if it is defined
    everywhere on it and can be traced without lifting the pencil. If you must lift the pencil, it is
    \emph{discontinuous} there, and the gap is called a \textbf{jump}.}
\end{vocabbox}

\vspace{0.10in}

\boxguard[12]
\begin{hookbox}
\small
\textbf{Two students, one paycheck, and \$60 that nobody can find.} A garden center pays \$12 an hour
for the first $40$ hours of a week, and \$18 an hour for every hour beyond $40$. A student worked
$45$ hours. \emph{What is the pay for that week?}

\smallskip
Student A: \textbf{``$45$ hours at \$12 is \$540, and there were $5$ overtime hours at \$18, which is
another \$90. So \$630.''}

\smallskip
Student B: \textbf{``$40$ hours at \$12 is \$480, and $5$ hours at \$18 is \$90. So \$570.''}

\smallskip
\textbf{Here is the question: which hours does each rate apply to?} Both students used both rates and
neither made an arithmetic mistake. Decide who is right, and say exactly which hours the other one
paid twice.

\ansline{Student B. Hours $41$--$45$ are overtime hours, so they are paid at \$18 \emph{instead of} \$12,}\\
\ansline{not on top of it. Student A paid those $5$ hours twice --- $5 \times \$12 = \$60$ too much.}\\
\end{hookbox}

\vspace{0.10in}

\boxguard[30]
\begin{notesbox}{1. One Function, Written in Pieces}
\small
Last night's delivery fee took two sentences to describe: \emph{flat \$5 for the first three miles,
then \$2 a mile after that.} There is a name for a function like that, and there is a way to write it
down.

\vspace{4pt}
A \textbf{piecewise-defined function} is \ans{one} function described by more than one rule, with
each rule applying only on its own stretch of the \ans{domain}. Each rule together with the stretch
of inputs it owns is called a \textbf{piece}.

\vspace{4pt}
\begin{center}
\begin{tikzpicture}
\begin{axis}[
    width=10.4cm, height=4.4cm,
    xlabel={$m$ (miles)}, ylabel={$F$ (fee, \$)},
    xmin=0, xmax=8, ymin=0, ymax=16,
    xtick={0,1,2,3,4,5,6,7,8}, ytick={0,5,10,15}, minor y tick num=4,
    grid=both, grid style={linegray!45}, minor grid style={linegray!30},
    axis lines=left,
    tick label style={font=\tiny}, label style={font=\scriptsize}]
  \addplot[royal, very thick] coordinates {(0,5) (3,5) (8,15)};
  \addplot[keyred, only marks, mark=*, mark size=1.7pt]
    coordinates {(1,5) (3,5) (4,7) (6,11) (8,15)};
\end{axis}
\end{tikzpicture}
\end{center}

\vspace{-4pt}
{\scriptsize Last night's item 10: $F(m)$ is the fee in dollars for a delivery of $m$ miles, for
deliveries up to $8$ miles.}

\vspace{4pt}
Here is the same graph written in symbols. Read the brace as the word \emph{if}:
\[
  F(m) = \begin{cases}
    5 & \text{if } 0 \le m \le 3,\\[2pt]
    5 + 2(m-3) & \text{if } 3 < m \le 8.
  \end{cases}
\]

\vspace{2pt}
\textbf{Which piece owns which inputs?} Fill the right-hand column.

\vspace{2pt}
\renewcommand{\arraystretch}{1.4}
\begin{tabularx}{\linewidth}{@{}p{4.4cm} X@{}}
\rowcolor{mist}
\textbf{Rule} & \textbf{The inputs it owns} \\
\hline
$F(m)=5$ & \ans{$0 \le m \le 3$, that is $[0,3]$} \\
$F(m)=5+2(m-3)$ & \ans{$3 < m \le 8$, that is $(3,8]$} \\
\end{tabularx}

\vspace{6pt}
\textbf{Evaluating takes one extra step, and it is always the same step.} Before you use a rule, ask
\textbf{which piece owns this input}. Then use that rule and \ans{only that one}.

\vspace{2pt}
\renewcommand{\arraystretch}{1.3}
\begin{tabular}{@{}|l|c|c|c|c|c|@{}}
\hline
$m$ (miles) & $1$ & $3$ & $4$ & $6$ & $8$ \\ \hline
Which piece? (1st or 2nd) & \ans{1st} & \ans{1st} & \ans{2nd} & \ans{2nd} &
    \ans{2nd} \\ \hline
$F(m)$ (\$) & \ans{$5$} & \ans{$5$} & \ans{$7$} & \ans{$11$} & \ans{$15$} \\
\hline
\end{tabular}

\vspace{6pt}
\textbf{The sentence that runs the whole lesson.} A piecewise function is \ans{one} function, not
several. The pieces are not allowed to overlap, so every input in the domain belongs to exactly
\ans{one} piece --- which means every input still has exactly \ans{one output}. That is Lesson 2.1's
definition of a function, and it has not changed.
\end{notesbox}

\vspace{0.10in}

\boxguard[30]
\begin{notesbox}{2. Where Each Rule Lives --- and Why the Paycheck Was Wrong}
\small
Here is the garden center's pay function. $P(h)$ is one week's pay in dollars for $h$ hours worked,
and the week is capped at $50$ hours.

\begin{center}
\begin{tikzpicture}
\begin{axis}[
    width=10.4cm, height=4.8cm,
    xlabel={$h$ (hours worked in the week)}, ylabel={$P$ (pay, \$)},
    xmin=0, xmax=50, ymin=0, ymax=700,
    xtick={0,5,10,15,20,25,30,35,40,45,50}, ytick={0,120,240,360,480,600},
    minor y tick num=1, minor x tick num=4,
    grid=both, grid style={linegray!45}, minor grid style={linegray!30},
    axis lines=left,
    tick label style={font=\tiny}, label style={font=\scriptsize}]
  \addplot[royal, very thick] coordinates {(0,0) (40,480) (50,660)};
  \addplot[keyred, only marks, mark=*, mark size=1.7pt]
    coordinates {(10,120) (20,240) (40,480) (45,570) (50,660)};
\end{axis}
\end{tikzpicture}
\end{center}

\vspace{-2pt}
\[
  P(h) = \begin{cases}
    12h & \text{if } 0 \le h \le 40,\\[2pt]
    480 + 18(h-40) & \text{if } 40 < h \le 50.
  \end{cases}
\]

\vspace{2pt}
\textbf{Read these five values off the graph.}

\vspace{2pt}
\renewcommand{\arraystretch}{1.3}
\begin{tabular}{@{}|l|c|c|c|c|c|@{}}
\hline
$h$ (hours) & $10$ & $20$ & $40$ & $45$ & $50$ \\ \hline
$P(h)$ (\$) & \ans{$120$} & \ans{$240$} & \ans{$480$} & \ans{$570$} & \ans{$660$} \\
\hline
\end{tabular}

\vspace{6pt}
\textbf{Now settle the hook.} Student A said \$630 for a $45$-hour week. The graph at $h=45$ says
\$\ans{$570$}. \textbf{Student \ans{B} answered the question.} Student A charged the \$12 rate
to \emph{all $45$} hours and then added overtime on top, so hours $41$ through $45$ got paid
\ans{twice}. \textbf{A rule applies only on the stretch of inputs it \ans{owns}} --- that is the
whole content of the word \emph{piecewise}, and it is the error to watch for all year.

\vspace{6pt}
\textbf{What happens at the boundary $h=40$?} Check both rules there: $12(40)=480$, and the second
rule starts at $480$ too. The two rules \ans{agree} at the boundary, so the graph has no break ---
you can trace it from end to end without lifting your pencil. A graph like that is called
\ans{continuous}.

\vspace{4pt}
\textbf{And here is the distinction to write down.} At $h=40$ the graph turns a sharp
\ans{corner} --- the line gets steeper because the pay rate went up. That is a change of
\emph{slope}, not a hole in the graph. \textbf{A corner is not a \ans{break}.}

\vspace{6pt}
\textbf{Every characteristic from this unit still works.} Fill the table --- one graph, all of Unit 2.

\vspace{2pt}
\renewcommand{\arraystretch}{1.4}
\begin{tabularx}{\linewidth}{@{}p{5.0cm} X@{}}
\rowcolor{mist}
\textbf{Characteristic} & \textbf{For $P$} \\
\hline
Domain (Lesson 2.1) & \ans{$0 \le h \le 50$, that is $[0,50]$} \\
Range (Lesson 2.1) & \ans{$0 \le P \le 660$, that is $[0,660]$} \\
Zeros and $y$-intercept (2.4) & \ans{one zero at $h=0$; $y$-intercept $(0,0)$ --- no hours, no pay} \\
Increasing / decreasing /\newline constant (2.4) & \ans{increasing on all of $[0,50]$; never constant
    or decreasing} \\
Absolute max and min (2.4) & \ans{max \$660 at $h=50$; min \$0 at $h=0$ --- both at endpoints} \\
End behavior, asymptotes (2.5) & \ans{none of either --- the domain stops at $50$, so nothing runs off} \\
\end{tabularx}
\end{notesbox}

\vspace{0.10in}

\boxguard[30]
\begin{notesbox}{3. When the Pieces Disagree --- Jumps, and What the Dots Mean}
\small
A shipping company charges by weight: \$6 for a package weighing up to and including $2$ pounds, \$10
for one over $2$ pounds up to and including $5$, and \$16 for one over $5$ pounds up to and including
$10$. Nothing in between --- the price changes all at once.

\begin{center}
\begin{tikzpicture}
\begin{axis}[
    width=10.4cm, height=4.6cm,
    xlabel={$w$ (pounds)}, ylabel={$S$ (shipping cost, \$)},
    xmin=0, xmax=10.6, ymin=0, ymax=20,
    xtick={0,1,2,3,4,5,6,7,8,9,10}, ytick={0,2,4,6,8,10,12,14,16,18},
    grid=both, grid style={linegray!45},
    axis lines=left,
    tick label style={font=\tiny}, label style={font=\scriptsize}]
  \addplot[royal, very thick] coordinates {(0,6) (2,6)};
  \addplot[royal, very thick] coordinates {(2,10) (5,10)};
  \addplot[royal, very thick] coordinates {(5,16) (10,16)};
  \addplot[royal, only marks, mark=*, mark size=2pt] coordinates {(2,6) (5,10) (10,16)};
  \addplot[royal, only marks, mark=*, mark size=2pt, mark options={fill=white}]
    coordinates {(0,6) (2,10) (5,16)};
\end{axis}
\end{tikzpicture}
\end{center}

\vspace{-4pt}
{\scriptsize A \textbf{filled} circle means that endpoint \emph{is} part of the graph. An
\textbf{open} circle means it is not.}

\vspace{4pt}
\textbf{The question the dots exist to answer: what is $S(2)$?} Look straight up from $w=2$. There
are two circles there --- one filled at \$6 and one open at \$10 --- so $S(2)=\$$\ans{$6$}. There
is exactly \ans{one} filled circle above every input, and there has to be, because a function
gives each input \ans{one output}.

\vspace{4pt}
\textbf{Fill in each cost.}

\vspace{2pt}
\renewcommand{\arraystretch}{1.3}
\begin{tabular}{@{}|l|c|c|c|c|c|c|@{}}
\hline
$w$ (pounds) & $1$ & $2$ & $3$ & $5$ & $7$ & $10$ \\ \hline
$S(w)$ (\$) & \ans{$6$} & \ans{$6$} & \ans{$10$} & \ans{$10$} & \ans{$16$} &
    \ans{$16$} \\
\hline
\end{tabular}

\vspace{6pt}
\textbf{Three things this graph does that last box's graph did not.}
\begin{enumerate}[leftmargin=1.5em, itemsep=3pt, topsep=3pt]
  \item \textbf{It breaks.} To trace it you must lift your pencil at $w=$ \ans{$2$} and at $w=$
        \ans{$5$}. A graph with a break in it is \ans{discontinuous}. The break itself is called a
        \textbf{jump}: here the cost jumps by \$4 and then by \$6.
  \item \textbf{Its range is a list, not an interval.} The only outputs this function ever produces
        are \ans{$\{6, 10, 16\}$}. No interval notation can say that --- you write the set, exactly as you did
        with the hot dogs in Lesson 2.1.
  \item \textbf{It is never increasing.} On each of its three stretches the graph is
        \ans{constant}, and a stretch is the only place ``increasing'' can be measured. The cost does
        go up --- but it goes up \emph{at} the jumps, and a jump happens over \ans{no interval} at all.
\end{enumerate}

\vspace{4pt}
\textbf{And the dots decide a bracket or a parenthesis.} The first piece owns the inputs
\ans{$0 < w \le 2$}: open at $0$ because a package with no weight is not a package, closed at $2$ because a
$2$-pound package pays \$6. That is Lesson 2.1's rule, back again --- a bracket \ans{includes}
the endpoint and a parenthesis leaves it out.
\end{notesbox}

\vspace{0.10in}

\boxguard[30]
\begin{practicebox}
\small
A delivery drone takes off, climbs, cruises at a fixed altitude, then descends to the ground. The
graph shows $A(t)$, its altitude in feet, $t$ seconds after take-off. It takes \emph{three} pieces.

\begin{center}
\begin{tikzpicture}
\begin{axis}[
    width=10.6cm, height=4.6cm,
    xlabel={$t$ (seconds after take-off)}, ylabel={$A$ (altitude, feet)},
    xmin=0, xmax=14, ymin=0, ymax=60,
    xtick={0,1,2,3,4,5,6,7,8,9,10,11,12,13,14}, ytick={0,10,20,30,40,50},
    minor y tick num=4,
    grid=both, grid style={linegray!45}, minor grid style={linegray!30},
    axis lines=left,
    tick label style={font=\tiny}, label style={font=\scriptsize}]
  \addplot[royal, very thick, domain=0:5, samples=60]{20*x-2*x^2};
  \addplot[royal, very thick] coordinates {(5,50) (9,50) (14,0)};
  \addplot[keyred, only marks, mark=*, mark size=1.7pt]
    coordinates {(0,0) (2,32) (5,50) (8,50) (11,30) (14,0)};
\end{axis}
\end{tikzpicture}
\end{center}

\begin{enumerate}[leftmargin=1.6em, itemsep=6pt, topsep=2pt]
  \item Read these altitudes off the graph, and name which piece --- climb, cruise, or descent ---
        owns each input.
\smallskip

\renewcommand{\arraystretch}{1.3}
\begin{tabular}{@{}|l|c|c|c|c|c|c|@{}}
\hline
$t$ (seconds) & $0$ & $2$ & $5$ & $8$ & $11$ & $14$ \\ \hline
$A(t)$ (feet) & \ans{$0$} & \ans{$32$} & \ans{$50$} & \ans{$50$} & \ans{$30$} &
    \ans{$0$} \\ \hline
Which piece? & \ans{climb} & \ans{climb} & \ans{climb} & \ans{cruise} & \ans{desc.} &
    \ans{desc.} \\
\hline
\end{tabular}

  \item Give the domain and the range of $A$, then give its zeros and its $y$-intercept. Say what
        each zero means about the flight.

\ansline{Domain $[0,14]$ seconds; range $[0,50]$ feet. Zeros at $t=0$ and $t=14$; $y$-intercept $(0,0)$.}\\
\ansline{The zero at $t=0$ is the moment of take-off and the zero at $t=14$ is the moment of landing ---}\\
\ansline{those are the two times the drone is on the ground.}\\

  \item Name the interval on which $A$ is increasing, the interval on which it is constant, and the
        interval on which it is decreasing. Then give the absolute maximum --- \emph{value and
        location}. Read the location carefully.

\ansline{Increasing on $[0,5]$, constant on $[5,9]$, decreasing on $[9,14]$.}\\
\ansline{Absolute maximum: $50$ feet. Its \emph{location} is not one point --- the drone is at $50$ feet}\\
\ansline{at every second from $t=5$ to $t=9$, so the maximum is attained on the whole interval $[5,9]$.}\\

  \item Is this graph continuous? Answer yes or no with a reason. Then answer the Lesson 2.5
        questions: what is the end behavior of $A$, and does it have any asymptotes?

\ansline{Yes --- the pieces meet at $t=5$ and $t=9$ with no gap, so you can trace the whole flight}\\
\ansline{without lifting your pencil. (There is a corner at $t=9$, and a corner is not a break.)}\\
\ansline{No end behavior and no asymptotes: the domain stops at $t=14$, so nothing runs off to infinity.}\\
\end{enumerate}
\end{practicebox}

\end{document}
